\section{Weighted empirical framework}
\label{sec:framework}

\subsection{Case weights and effective sample size}

Let $(X_i,Y_i)$, $i=1,\ldots,n$, denote the observed pairs and let
$w_i\geq 0$ be their case weights, with at least one positive weight. All
statistics below depend on the weights only through the normalized masses
$p_i = w_i/\sum_{j=1}^n w_j$, which define the weighted empirical measure
$\mathbb{P}_{n,w} = \sum_{i=1}^n p_i\,\delta_{(X_i,Y_i)}$. They are therefore
invariant under multiplication of all weights by a common positive constant,
and an observation with zero weight is treated as absent. The weights are case
masses, not weights attached to rank positions; the latter are used by
coefficients that deliberately emphasize the top or bottom of a ranking
\citep{shieh1998} and are not considered here.

The concentration of the masses is summarized by Kish's effective sample size
\citep{kish1965},
\begin{align*}
    n_{\mathrm{eff}}
    &= \frac{\bigl(\sum_{i=1}^n w_i\bigr)^2}{\sum_{i=1}^n w_i^2}
    = \frac{1}{\sum_{i=1}^n p_i^2}.
\end{align*}
It equals $n$ under uniform weights and decreases as mass concentrates on
fewer observations. Its role is explained by the weighted mean: if the
observations are independent and identically distributed, the weights are
fixed, and $h(X_i,Y_i)$ has variance $\sigma_h^2$, then
\begin{align*}
    \operatorname{Var}\bigl(\mathbb{P}_{n,w}h \mid p_1,\ldots,p_n\bigr)
    &= \sigma_h^2\sum_{i=1}^n p_i^2
    = \frac{\sigma_h^2}{n_{\mathrm{eff}}}.
\end{align*}
Thus $n_{\mathrm{eff}}$ is the exact precision-equivalent sample size for a
weighted mean. Section~\ref{sec:independence-tests} identifies when this
scalar also gives the correct normalization for a dependence statistic and
when further weight functionals are required.

\subsection{Weighted ranks}
\label{subsec:weighted-ranks}

The strict and weak weighted empirical distribution functions of $X$ are
\begin{align*}
    F_{X,w}^{-}(x)
    &= \sum_{j=1}^n p_j\,\mathbf{1}(X_j<x),
    &
    F_{X,w}(x)
    &= \sum_{j=1}^n p_j\,\mathbf{1}(X_j\leq x),
\end{align*}
and analogously for $Y$. The strict weighted bivariate rank at observation
$i$,
\begin{align*}
    H_{i,w}^{-}
    &= \sum_{j=1}^n p_j\,\mathbf{1}(X_j<X_i,\,Y_j<Y_i),
\end{align*}
is the basic computational quantity of the weighted Hoeffding statistic in
Section~\ref{sec:weighted-hoeffding}.

Spearman's coefficient and Blomqvist's beta use a weighted average-rank
convention for ties. Let $G_i^X=\{j\colon X_j=X_i\}$ be the tie group
containing $X_i$, and let
\begin{align*}
    P_i^X
    &= \sum_{j\in G_i^X}p_j,
    &
    E_i^X
    &= \sum_{\substack{j<k\\ j,k\in G_i^X}}p_jp_k
\end{align*}
be its total mass and pair mass. Every member of the tie group receives the
zero-based average score
\begin{align*}
    R_i^X
    &= F_{X,w}^{-}(X_i)+\frac{E_i^X}{P_i^X},
\end{align*}
and $R_i^Y$ is defined in the same way. Under uniform weights, $nR_i^X$ is the
usual zero-based midrank. With unequal weights, the pair-mass term is part of
the definition and should not be replaced by one half of the mass of the tie
group.

Algorithm~\ref{alg:weighted-ranks} computes the strict ranks, weak ranks, and
average scores in one scan over tie groups after sorting. For a tie group $G$,
its total mass and pair mass are denoted by $P_G$ and $E_G$, and $P_{<G}$ is
the mass of all preceding groups. The running time is $O(n\log n)$ and the
auxiliary storage is $O(n)$. The same scan computes strict rank sums for other
masses, such as $p_i^a$, by replacing $p_i$ by the desired mass; normalization
matters only when the output is interpreted as an empirical distribution
function.

\begin{algorithm}[ht]
  \caption{Weighted marginal ranks and average scores}
  \label{alg:weighted-ranks}
  \small
  \begin{algorithmic}[1]
    \Procedure{WeightedRanks}{$(X_i,p_i)_{i=1}^n$}
      \State Sort the records by increasing $X_i$
      \State $P_{<G}\gets0$
      \ForAll{maximal tie groups $G$, in increasing order}
        \State $P_G\gets\sum_{j\in G}p_j$
        \State $E_G\gets\bigl(P_G^2-\sum_{j\in G}p_j^2\bigr)/2$
        \ForAll{$i\in G$}
          \State $F_{X,w}^{-}(X_i)\gets P_{<G}$
          \State $F_{X,w}(X_i)\gets P_{<G}+P_G$
          \State $R_i^X\gets P_{<G}+E_G/P_G$
        \EndFor
        \State $P_{<G}\gets P_{<G}+P_G$
      \EndFor
      \State \Return the three score vectors in the original observation order
    \EndProcedure
  \end{algorithmic}
\end{algorithm}

\subsection{Distinct-index sums}
\label{subsec:distinct-index-sums}

Weighted statistics of order two or higher involve sums over distinct
observations. For $k\leq n$, the elementary symmetric sum
\begin{align*}
    e_k(p_1,\ldots,p_n)
    &= \sum_{1\leq i_1<\cdots<i_k\leq n} p_{i_1}\cdots p_{i_k}
\end{align*}
is the total mass of unordered $k$-tuples of distinct observations, and
$k!e_k(p_1,\ldots,p_n)$ is the mass of ordered tuples. In particular,
$e_2(p_1,\ldots,p_n) = (1-\sum_{i=1}^n p_i^2)/2$. These sums normalize the
weighted Kendall and Hoeffding statistics;
Section~\ref{sec:weighted-hoeffding} evaluates them from power sums of the
masses by Newton's identities.
